\textbf{COLUNI 2025 -- Questão 33}\par\medskip Onde começa um rio? De onde vem a água que o preenche? Para onde ele segue descendo? Da água que cai sobre a terra, uma parte escorre pela superfície, outra infiltra no solo e percorre um caminho subterrâneo. Em algum lugar, a parte da água infiltrada poderá aparecer novamente: a nascente é o surgimento de um rio.

\par

Analise as afirmativas a respeito deste importante elemento natural e assinale a INCORRETA:\begin{enumerate}
\item À medida que o curso d’água segue seu fluxo, indo da mina até desaguar na foz, sua vazão aumenta no
      leito e sua altitude diminui.
\item As características das correntes fluviais podem variar: volume, cor, temperatura, biota, profundidade,
      navegabilidade, etc.
\item Ribeiros, córregos e riachos são nomes utilizados para denominar os precursores e afluentes de rios
      caudalosos.
\item O regime fluvial diferencia os rios em permanentes e perenes, permitindo também distinguir as fases do
      seu ciclo entre cheias e vazantes.\end{enumerate}
